\documentclass[10pt,a4paper]{amsart}
\usepackage{amsmath,amssymb,amsthm,graphicx}
\usepackage{verbatim}
\newenvironment{reference}{\comment}{\endcomment}
%\newenvironment{reference}{}{}
\newenvironment{slogan}{\comment}{\endcomment}
\newenvironment{history}{\comment}{\endcomment}

% For commutative diagrams we use Xy-pic
\usepackage[all]{xy}

% We use 2cell for 2-commutative diagrams.
\xyoption{2cell}
% Theorem environments.
%
\theoremstyle{plain}
\newtheorem{theorem}[subsection]{Theorem}
\newtheorem{proposition}[subsection]{Proposition}
\newtheorem{lemma}[subsection]{Lemma}

\theoremstyle{definition}
\newtheorem{definition}[subsection]{Definition}
\newtheorem{example}[subsection]{Example}
\newtheorem{exercise}[subsection]{Exercise}
\newtheorem{situation}[subsection]{Situation}

\theoremstyle{remark}
\newtheorem{remark}[subsection]{Remark}
\newtheorem{remarks}[subsection]{Remarks}

\numberwithin{equation}{subsection}

% Macros
%
\def\lim{\mathop{\mathrm{lim}}\nolimits}
\def\colim{\mathop{\mathrm{colim}}\nolimits}
\def\Spec{\mathop{\mathrm{Spec}}}
\def\Hom{\mathop{\mathrm{Hom}}\nolimits}
\def\Ext{\mathop{\mathrm{Ext}}\nolimits}
\def\SheafHom{\mathop{\mathcal{H}\!\mathit{om}}\nolimits}
\def\SheafExt{\mathop{\mathcal{E}\!\mathit{xt}}\nolimits}
\def\Sch{\mathit{Sch}}
\def\Mor{\mathop{\mathrm{Mor}}\nolimits}
\def\Ob{\mathop{\mathrm{Ob}}\nolimits}
\def\Sh{\mathop{\mathit{Sh}}\nolimits}
\def\NL{\mathop{N\!L}\nolimits}
\def\CH{\mathop{\mathrm{CH}}\nolimits}
\def\proetale{{pro\text{-}\acute{e}tale}}
\def\etale{{\acute{e}tale}}
\def\QCoh{\mathit{QCoh}}
\def\Ker{\mathop{\mathrm{Ker}}}
\def\Im{\mathop{\mathrm{Im}}}
\def\Coker{\mathop{\mathrm{Coker}}}
\def\Coim{\mathop{\mathrm{Coim}}}

% Boxtimes
%
\DeclareMathSymbol{\boxtimes}{\mathbin}{AMSa}{"02}

%
% Macros for moduli stacks/spaces
%
\def\QCohstack{\mathcal{QC}\!\mathit{oh}}
\def\Cohstack{\mathcal{C}\!\mathit{oh}}
\def\Spacesstack{\mathcal{S}\!\mathit{paces}}
\def\Quotfunctor{\mathrm{Quot}}
\def\Hilbfunctor{\mathrm{Hilb}}
\def\Curvesstack{\mathcal{C}\!\mathit{urves}}
\def\Polarizedstack{\mathcal{P}\!\mathit{olarized}}
\def\Complexesstack{\mathcal{C}\!\mathit{omplexes}}
% \Pic is the operator that assigns to X its picard group, usage \Pic(X)
% \Picardstack_{X/B} denotes the Picard stack of X over B
% \Picardfunctor_{X/B} denotes the Picard functor of X over B
\def\Pic{\mathop{\mathrm{Pic}}\nolimits}
\def\Picardstack{\mathcal{P}\!\mathit{ic}}
\def\Picardfunctor{\mathrm{Pic}}
\def\Deformationcategory{\mathcal{D}\!\mathit{ef}}
\usepackage{cleveref}
\crefname{originallemma}{original lemma}{original lemmas}
\Crefname{originallemma}{Original lemma}{Original lemmas}
\makeatletter
\newlabel{lemma-dimension-at-a-point-finite-type-over-field}{{00OT}{source}{}{sourcepointer-1}{}}
\newlabel{lemma-dimension-at-a-point-finite-type-over-field@cref}{{[originallemma][0][]00OT}{[1][0][]source}{}{}{}}
\makeatother
\crefname{originallemma}{original lemma}{original lemmas}
\Crefname{originallemma}{Original lemma}{Original lemmas}
\makeatletter
\newlabel{lemma-dimension-prime-polynomial-ring}{{00P0}{source}{}{sourcepointer-2}{}}
\newlabel{lemma-dimension-prime-polynomial-ring@cref}{{[originallemma][0][]00P0}{[1][0][]source}{}{}{}}
\makeatother
\crefname{originallemma}{original lemma}{original lemmas}
\Crefname{originallemma}{Original lemma}{Original lemmas}
\makeatletter
\newlabel{lemma-dimension-spell-it-out}{{00OS}{source}{}{sourcepointer-3}{}}
\newlabel{lemma-dimension-spell-it-out@cref}{{[originallemma][0][]00OS}{[1][0][]source}{}{}{}}
\makeatother
\title{The local dimension at a point is the dimension of the local ring plus the transcendence degree of its residue field}
\author{The Stacks Project Authors}
\date{stacks-project:algebra.tex:local-snapshot-20261012; GNU FDL 1.2 or later}
\begin{document}
\maketitle
\noindent\textit{Editorial note.} All delivered mathematical text is the author's own, copied from the pinned Stacks Project chapter file algebra.tex; no statement or proof step was rewritten, reordered or removed.\par\medskip
\noindent\textit{Editorial note.} Preamble provenance, as required for this source: the frozen snapshot stages the chapter's own preamble file beside it (preamble.tex, 217 lines, 8116 bytes, the same shared Stacks-wide preamble the published chow and etale-cohomology cards use); the first line of the pinned chapter file is \textbackslash{}input\{preamble\}, so no preamble text lives inside the chapter itself. The card takes that file's lines 8-18 (the verbatim package and the three \textbackslash{}comment-based environments reference/slogan/history at lines 9-12, section 13.2, together with the Xy-pic setup) and its lines 154 to the runtime end of the file (the native theorem declarations and mathematical macros). Not one character was changed, no package was added, no macro was defined, and the excerpt serves only as preamble, never as a span coordinate.\par\medskip
\noindent\textit{Editorial note.} The same-chapter pointers are located at their \textbackslash{}label lines inside the pinned chapter file; the cross-chapter pointers are located in the project's own tag list staged in the snapshot (sources/stacks-chow/tags.txt), at the verbatim lines carrying their canonical tags.\par\medskip
\noindent\textit{Editorial note.} The counted proof refers to two same-chapter statements by label; both are recorded as navigation stubs pointed at their \textbackslash{}label lines, and the arguments consuming them stay inside the counted span.\par\medskip
\noindent\textit{Editorial note.} The counted spans contain no \textbackslash{}cite, so the card carries no bibliography.\par\medskip
\noindent\textit{Allowed standard prerequisites (editorial).}\begin{itemize}
\item The dimension of the closed set defined by a prime as a transcendence degree. Invoked at the start of the counted proof (lines 28452-28454). Identifies the transcendence degree of the residue field with the dimension of the vanishing set of the prime.
\item Maximal chains of primes above a given prime. Used in the counted proof (lines 28455-28478). Estimates the local dimension by comparing chains in the spectrum with chains in the local ring.
\end{itemize}\medskip
\ifdefined\hypertarget\hypertarget{sourcepointer-1}{}\fi
\noindent\textit{Source reference.} Original lemma 00OT: sources/stacks-algebra/pinned-original-source.tex, source lines 28004--28004. Navigation pointer for \textbackslash{}ref\{lemma-dimension-at-a-point-finite-type-over-field\} used inside the counted span (line 28463); the argument that consumes it is retained in the span.\par\medskip
\ifdefined\hypertarget\hypertarget{sourcepointer-2}{}\fi
\noindent\textit{Source reference.} Original lemma 00P0: sources/stacks-algebra/pinned-original-source.tex, source lines 28395--28395. Navigation pointer for \textbackslash{}ref\{lemma-dimension-prime-polynomial-ring\} used inside the counted span (line 28452); the argument that consumes it is retained in the span.\par\medskip
\ifdefined\hypertarget\hypertarget{sourcepointer-3}{}\fi
\noindent\textit{Source reference.} Original lemma 00OS: sources/stacks-algebra/pinned-original-source.tex, source lines 27980--27980. Navigation pointer for \textbackslash{}ref\{lemma-dimension-spell-it-out\} used inside the counted span (line 28458); the argument that consumes it is retained in the span.\par\medskip
\begin{lemma}
\label{lemma-dimension-at-a-point-finite-type-field}
Let $k$ be a field.
Let $S$ be a finite type $k$ algebra.
Let $X = \Spec(S)$.
Let $\mathfrak p \subset S$ be a prime ideal,
and let $x \in X$ be the corresponding point.
Then we have
$$
\dim_x(X) = \dim(S_{\mathfrak p}) + \text{trdeg}_k\ \kappa(\mathfrak p).
$$
\end{lemma}

\begin{proof}
By Lemma \ref{lemma-dimension-prime-polynomial-ring} we know that
$r = \text{trdeg}_k\ \kappa(\mathfrak p)$ is equal to the
dimension of $V(\mathfrak p)$.
Pick any maximal chain of primes
$\mathfrak p \subset \mathfrak p_1 \subset \ldots \subset \mathfrak p_r$
starting with $\mathfrak p$ in $S$.
This has length $r$ by Lemma \ref{lemma-dimension-spell-it-out}.
Let $\mathfrak q_j$, $j \in J$ be the minimal
primes of $S$ which are contained in $\mathfrak p$.
These correspond $1-1$ to minimal primes in $S_{\mathfrak p}$
via the rule $\mathfrak q_j \mapsto \mathfrak q_jS_{\mathfrak p}$.
By Lemma \ref{lemma-dimension-at-a-point-finite-type-over-field}
we know that $\dim_x(X)$ is equal
to the maximum of the dimensions of the rings $S/\mathfrak q_j$.
For each $j$ pick a maximal chain of primes
$\mathfrak q_j \subset \mathfrak p'_1 \subset \ldots \subset \mathfrak p'_{s(j)}
= \mathfrak p$.
Then $\dim(S_{\mathfrak p}) = \max_{j \in J} s(j)$.
Now, each chain
$$
\mathfrak q_j \subset \mathfrak p'_1 \subset \ldots \subset
\mathfrak p'_{s(j)} = \mathfrak p \subset
\mathfrak p_1 \subset \ldots \subset \mathfrak p_r
$$
is a maximal chain in $S/\mathfrak q_j$, and by what was said
before we have
$\dim_x(X) = \max_{j \in J} r + s(j)$.
The lemma follows.
\end{proof}

\end{document}
